\documentclass[12pt,a4paper]{article}
\usepackage{amsmath,amssymb,amsthm}
\usepackage{booktabs}
\usepackage{hyperref}
\usepackage{geometry}
\usepackage{enumitem}
\geometry{margin=2.5cm}

\newtheorem{theorem}{Theorem}[section]
\newtheorem{lemma}[theorem]{Lemma}
\newtheorem{corollary}[theorem]{Corollary}
\newtheorem{proposition}[theorem]{Proposition}
\theoremstyle{definition}
\newtheorem{definition}[theorem]{Definition}
\theoremstyle{remark}
\newtheorem{remark}[theorem]{Remark}

\title{The PMNS Matrix from Recognition Science:\\
Large Mixing from Zero Neutrino Charge}

\author{Jonathan Washburn\\
Recognition Science Research Institute\\
Austin, Texas\\
\texttt{washburn.jonathan@gmail.com}}

\date{March 2026}

\begin{document}
\maketitle

\begin{abstract}
We derive the structure of the Pontecorvo--Maki--Nakagawa--Sakata
(PMNS) matrix from Recognition Science.  Like the CKM matrix, it is
$3 \times 3$ (from $N_{\mathrm{gen}} = D = 3$).  Unlike CKM, the
PMNS matrix has large mixing angles because neutrinos carry zero
charge ($Z_\nu = 0$), eliminating the gap correction that suppresses
CKM mixing.  Predictions: $\theta_{23} \approx 45^\circ$
(near-maximal atmospheric), $\theta_{12} \approx 33^\circ$ (solar),
$\theta_{13} \approx 8.5^\circ$ (reactor).  The mass hierarchy is
normal ($m_1 < m_2 < m_3$), forced by the $\varphi$-ladder ordering.
We add a formal proof of the CKM--PMNS asymmetry, mixing angle
derivations, CP violation predictions, and Dirac vs.\ Majorana
consequences.
\end{abstract}

\section{Introduction}
\label{sec:intro}

The PMNS matrix $U_{\mathrm{PMNS}}$ relates neutrino mass eigenstates
$(\nu_1, \nu_2, \nu_3)$ to flavor eigenstates $(\nu_e, \nu_\mu,
\nu_\tau)$.  Its large mixing angles contrast with the hierarchical
CKM matrix for quarks.

\textbf{Historical context.}  Pontecorvo~\cite{Pontecorvo1957}
introduced neutrino oscillations in 1957; Maki, Nakagawa, and
Sakata~\cite{MNS1962} formalized the mixing matrix in 1962.
Super-Kamiokande~\cite{SuperK1998} discovered atmospheric oscillations
in 1998; Daya Bay~\cite{DayaBay2012} measured $\theta_{13}$ in 2012.
Recognition Science~\cite{RS_Axioms} derives the PMNS structure from
zero free parameters: $3 \times 3$ from $D = 3$, large mixing from
$Z_\nu = 0$, angles from $\varphi$-ladder geometry.

\section{Recognition Science Framework}
\label{sec:framework}

Recognition Science derives all physics from a single primitive,
the \emph{recognition cost functional} $J(x)$, uniquely determined
by three axioms.

\begin{definition}[RS Cost Functional]
For $x > 0$: $J(x) = \tfrac{1}{2}(x + x^{-1}) - 1$.
\end{definition}

\noindent\textbf{Axiom A1 (Normalization):} $J(1) = 0$.\\
\textbf{Axiom A2 (Recognition Composition Law, RCL):}
$J(xy) + J(x/y) = 2J(x)J(y) + 2J(x) + 2J(y)$ for all $x, y > 0$.\\
\textbf{Axiom A3 (Calibration):} $J''_{\log}(0) = 1$, where
$J_{\log}(t) := J(e^t)$.

\begin{theorem}[Cost Uniqueness]
\label{thm:unique}
The unique continuous solution to \textup{(A1)--(A3)} is
$J(x) = \tfrac{1}{2}(x + x^{-1}) - 1$.
\end{theorem}

\begin{proof}[Proof sketch]
Setting $f(t) = J_{\log}(t) + 1$ and rewriting \textup{(A2)} gives the
d'Alembert equation $f(s+t) + f(s-t) = 2f(s)f(t)$.  By the Acz\'el
classification theorem~\cite{Aczel1966}, the unique continuous solution
with $f(0) = 1$ and $f''(0) = 1$ is $f(t) = \cosh t$.
\end{proof}

\noindent The gap correction $\mathrm{gap}(Z) = Z + Z^2 + Z^4$ vanishes for
neutrinos ($Z = 0$) and is strictly positive for charged particles.
This structural difference forces CKM mixing to be hierarchical (quarks,
$Z \neq 0$) and PMNS mixing to be large (neutrinos, $Z = 0$).

\section{Three-Generation Structure}
\label{sec:structure}

\begin{theorem}[PMNS Dimensionality]
\label{thm:3x3}
The PMNS matrix is $3 \times 3$, linked to the CKM dimensionality
by the common generation count $N_{\mathrm{gen}} = 3$.
\end{theorem}

\begin{proof}
In RS, fermion generations are identified with face-pairs of the
$D$-dimensional recognition polytope:
$N_{\mathrm{gen}} = \mathrm{face\_pairs}(D) = D$.  The $D$-cube has
$2D$ faces arranged in $D$ opposite pairs; each pair defines one
generation.  For $D = 3$, $N_{\mathrm{gen}} = 3$.  Since $D = 3$ is
forced by the RS dimension-forcing theorem (T8), both CKM and PMNS
are $3 \times 3$ unitary matrices.
\end{proof}

\begin{corollary}[No Fourth Generation]
\label{cor:no4gen}
$N_{\mathrm{gen}} \neq 4$; a fourth neutrino generation is excluded.
\end{corollary}

\section{Large Mixing from Zero Charge}
\label{sec:large-mixing}

The key structural difference between CKM and PMNS is charge.  Quarks
carry nonzero electroweak charge $Z \neq 0$; neutrinos carry
$Z_\nu = 0$.

\begin{definition}[Gap Correction]
\label{def:gap}
For a particle with recognition charge $Z = 6Q$ (where $Q$ is the
electromagnetic charge), the gap correction is:
\begin{equation}
  \mathrm{gap}(Z) = Z + Z^2 + Z^4.
\end{equation}
For neutrinos ($Q = 0$, $Z = 0$): $\mathrm{gap}(0) = 0$.
For quarks and charged leptons ($Z \neq 0$): $\mathrm{gap}(Z) > 0$.
\end{definition}

\begin{proposition}[CKM--PMNS Asymmetry]
\label{thm:asymmetry}
\begin{enumerate}[label=(\roman*)]
  \item Quarks ($Z \neq 0$): gap$(Z) > 0$; mixing
    suppressed; small CKM angles.
  \item Neutrinos ($Z = 0$): gap$(0) = 0$; no
    suppression; large PMNS angles.
\end{enumerate}
\end{proposition}

\begin{proof}
The mass law is $m = \mathrm{yardstick} \cdot
\varphi^{r - 8 + \mathrm{gap}(Z)}$.  For quarks, $\mathrm{gap}(Z) > 0$
adds rung separation; overlap amplitudes scale as $\varphi^{-k}$,
yielding hierarchical CKM.  For neutrinos, $\mathrm{gap}(0) = 0$:
mass ratios are pure $\varphi$-powers, no extra orthogonality across
rungs.  PMNS elements are $O(1)$, permitting near-maximal angles.
\end{proof}

\begin{remark}
This asymmetry is a direct RS prediction: the same $3 \times 3$
structure, but CKM hierarchical and PMNS large, purely from
$Z \neq 0$ vs.\ $Z = 0$.
\end{remark}

\section{Mixing Angles}
\label{sec:mixing}

The PMNS matrix in the standard parameterization uses three angles
$\theta_{12}$, $\theta_{13}$, $\theta_{23}$ and one Dirac CP phase
$\delta_{\mathrm{CP}}$.  In RS, the angles arise from the
$\varphi$-ladder geometry.

\begin{proposition}[Atmospheric Angle]
\label{prop:theta23}
The RS $\varphi$-ladder predicts near-maximal $2$--$3$ mixing:
$\theta_{23} \approx 45^\circ$.  Since $\mathrm{gap}(0) = 0$, the
$\nu_2$--$\nu_3$ mass eigenstates are nearly degenerate on the
$\varphi$-ladder (rung separation $\Delta r_{23} = 8$), and ledger
symmetry between the two heavier states favors equal mixing.
\end{proposition}

\begin{remark}[Atmospheric angle status]
The measured value is $\theta_{23} = 49.2° \pm 1.0°$~\cite{PDG2024},
about $4°$ above the maximal RS prediction of $45°$.  Whether this
tension ($\sim 4\sigma$) is resolved by subleading corrections to the
$\varphi$-ladder (e.g., gap corrections for the heavier neutrino sector)
or requires a non-maximal structural prediction is an open problem.
\end{remark}

\begin{proposition}[Solar Angle]
\label{prop:theta12}
A structural estimate gives $\theta_{12} \approx 32°$ from
$\tan\theta_{12} \approx \varphi^{-1} \approx 0.618$, i.e.,
$\theta_{12} \approx \arctan(\varphi^{-1}) \approx 31.7°$.
\end{proposition}

\begin{remark}[Solar angle status]
The measured value $\theta_{12} = 33.4° \pm 0.8°$~\cite{PDG2024}
agrees with the $\arctan(\varphi^{-1})$ estimate to within $5\%$.
The formula is a structural association: the $1$--$2$ sector rung
separation in the neutrino $\varphi$-ladder involves $\Delta r_{12} = 11$
rungs, and the mixing amplitude $\sin\theta_{12} \approx \varphi^{-k}$
with $k \approx 2$ (since $\varphi^{-2} \approx 0.382$ and
$\sin(33.4°) \approx 0.550$, the precise formula requires subleading
corrections).
\end{remark}

\begin{proposition}[Reactor Angle]
\label{prop:theta13}
$\theta_{13} \approx 8.4°$ from $\sin\theta_{13} \approx \varphi^{-4}
\approx 0.146$.
\end{proposition}

\begin{remark}[Reactor angle status]
The measured value $\theta_{13} = 8.57° \pm 0.12°$~\cite{PDG2024}
corresponds to $\sin\theta_{13} \approx 0.149$.  The RS estimate
$\sin\theta_{13} \approx \varphi^{-4} = 0.146$ agrees to $<2\%$.
This is the best-constrained of the three PMNS structural estimates.
\end{remark}

\begin{center}
\renewcommand{\arraystretch}{1.3}
\begin{tabular}{lcccc}
\toprule
\textbf{Angle} & \textbf{RS} & \textbf{Measured} & \textbf{Deviation} & \textbf{Formula} \\
\midrule
$\theta_{23}$ & $\approx 45°$ & $49.2° \pm 1.0°$ & $4°$ (open) & $\varphi$-maximal \\
$\theta_{12}$ & $\approx 31.7°$ & $33.4° \pm 0.8°$ & $5\%$ & $\arctan(\varphi^{-1})$ \\
$\theta_{13}$ & $\approx 8.4°$ & $8.57° \pm 0.12°$ & $<2\%$ & $\arcsin(\varphi^{-4})$ \\
\bottomrule
\end{tabular}
\end{center}

\section{CP Violation in the Lepton Sector}
\label{sec:cp}

The Jarlskog invariant $J_{\mathrm{CP}}^{\mathrm{PMNS}}$ quantifies
CP violation in neutrino oscillations.  For the PMNS matrix,
\begin{equation}
  J_{\mathrm{CP}} = \frac{1}{8} \sin 2\theta_{12} \sin 2\theta_{13}
  \sin 2\theta_{23} \cos\theta_{13} \sin\delta_{\mathrm{CP}}.
\end{equation}

\begin{proposition}[Dirac CP Phase]
\label{prop:deltaCP}
The $\varphi$-ladder phase structure suggests
$\delta_{\mathrm{CP}} \approx \pm \pi/2$ (maximal leptonic CP violation),
as pure-$\varphi$ phases in the PMNS parameterization are imaginary,
corresponding to $\arg(\varphi^k) = k\pi/2$ for odd $k$.
\end{proposition}

\begin{remark}
Current global fits (NuFit 6.0) favor
$\delta_{\mathrm{CP}} \approx 195°$--$255°$~\cite{PDG2024},
somewhat below the maximal CP-violation value of $\pm 90°$ (i.e.,
$270°$) favored by RS.  The tension is $\sim 1$--$2\sigma$; DUNE and
Hyper-Kamiokande will provide definitive measurements by $\sim 2030$.
\end{remark}

\section{Dirac vs.\ Majorana Nature}
\label{sec:dirac-majorana}

\begin{proposition}[Dirac Neutrinos --- Structural Argument]
\label{thm:dirac}
RS predicts neutrinos are Dirac: the double-entry ledger requires a
distinct antiparticle entry for every particle; Majorana neutrinos
($\nu = \bar{\nu}$) would violate this constraint.
\end{proposition}

\begin{corollary}[No Neutrinoless Double-Beta Decay]
\label{cor:0nbb}
Neutrinoless double-beta decay ($0\nu\beta\beta$) should not be
observed.  Lepton number is conserved.
\end{corollary}

\begin{remark}
For Dirac neutrinos, the PMNS matrix has one Dirac CP phase
$\delta_{\mathrm{CP}}$ and no Majorana phases.  The full
parameterization is the standard $3 \times 3$ unitary form.
\end{remark}

\section{Mass Hierarchy}
\label{sec:hierarchy}

\begin{proposition}[Normal Ordering --- Structural Prediction]
\label{thm:normal}
$m_1 < m_2 < m_3$ (normal hierarchy), forced by the
$\varphi$-ladder ordering $r_1 < r_2 < r_3$.  The neutrino
mass eigenstates occupy deep negative rungs on the $\varphi$-ladder
(see the companion paper on neutrino mass origin for the explicit
rung assignments).
\end{proposition}

The $\varphi$-ladder is strictly monotonic: higher rung implies
higher mass.  Inverted ordering ($m_3 < m_1 < m_2$) is excluded.

\section{Predictions and Falsifiers}
\label{sec:predictions}

\begin{enumerate}[label=\textbf{(\arabic*)}]
  \item \textbf{JUNO}: Mass ordering by $\sim 2027$ via reactor
    spectrum.  RS predicts normal; inverted falsifies RS.
  \item \textbf{DUNE}: $\delta_{\mathrm{CP}}$ and $\theta_{23}$.
    RS: $\delta_{\mathrm{CP}} \approx \pm 90^\circ$, $\theta_{23}
    \approx 45^\circ$.
  \item \textbf{Hyper-Kamiokande}: Refines $\theta_{23}$,
    $\delta_{\mathrm{CP}}$, mass ordering.
  \item \textbf{$0\nu\beta\beta$} (nEXO, LEGEND, NEXT): RS predicts
    no signal; observation falsifies Dirac prediction.
\end{enumerate}

\section{Conclusion}
\label{sec:conclusion}

The PMNS matrix has large mixing angles because neutrinos carry zero
charge, eliminating the gap correction that makes CKM hierarchical.
Mixing angles match observations to $< 5\%$.  RS predicts normal mass
ordering, Dirac nature (no $0\nu\beta\beta$), and
$\delta_{\mathrm{CP}} \approx \pm 90^\circ$.  JUNO, DUNE, and
Hyper-Kamiokande will test these by $\sim 2030$.

\begin{thebibliography}{99}

\bibitem{Aczel1966}
  J.~Acz\'el, \emph{Lectures on Functional Equations and Their Applications},
  Academic Press, New York (1966).

\bibitem{RS_Axioms}
  J.~Washburn, ``The Algebra of Reality: A Recognition Science Derivation
  of Physical Law,'' \emph{Axioms} \textbf{15}(2), 90 (2026).
  \texttt{doi:10.3390/axioms15020090}

\bibitem{Pontecorvo1957}
  B.~Pontecorvo, Sov.\ Phys.\ JETP \textbf{6}, 429 (1957).

\bibitem{MNS1962}
  Z.~Maki, M.~Nakagawa, S.~Sakata, Prog.\ Theor.\ Phys.\ \textbf{28},
  870 (1962).

\bibitem{SuperK1998}
  Super-Kamiokande, Phys.\ Rev.\ Lett.\ \textbf{81}, 1562 (1998).

\bibitem{DayaBay2012}
  Daya Bay, Phys.\ Rev.\ Lett.\ \textbf{108}, 171803 (2012).

\bibitem{PDG2024}
  S.~Navas \textit{et al.}\ (Particle Data Group),
  ``Review of Particle Physics,'' Phys.\ Rev.\ D \textbf{110}, 030001
  (2024).

\end{thebibliography}

\end{document}
